\subsection{Sprint 3}

Com tudo o que aconteceu no Rio Grande do Sul, o time foi afetado de diferentes
formas, e muitos colegas tinham dificuldades para participar das aulas online. A nossa
\ac{ages} IV, Sofia, entrou em contato individualmente com cada membro perguntando
como estávamos e se conseguiríamos participar normalmente. Apesar de eu não ter sido
afetado diretamente pelas enchentes, estava com parentes temporariamente em casa,
pois tinham perdido tudo --- casa e loja --- na enchente em Mathias Velho, Canoas.

Confirmei minha presença e que trabalharia o máximo possível. Logo após a
retrospectiva, o time dividiu as tasks e fiquei responsável por implementar a
\textbf{Issue 17 --- Solicitação de Falta}.

Comecei a desenvolver uma estratégia para implementar a funcionalidade e decidi criar
uma nova entidade: \textit{absenceNotice}. Ao apresentar a proposta ao Pedro, ele
sugeriu uma abordagem diferente e mais enxuta: \textbf{reutilizar e completar a
entidade \textit{Request}}, que estava incompleta no projeto. Segui a sugestão e,
de fato, economizei muito código, além de aproveitar uma estrutura que já fazia parte
do sistema.

Ao finalizar o desenvolvimento, precisava realizar os testes, mas me deparei com um
problema: os POSTs não estavam funcionando. Após investigar, descobri que precisava
\textbf{derrubar e reinstanciar todo o Docker}. Após fazer isso, o erro persistiu ---
dessa vez, porque havia sido implementada autenticação e era necessário cadastrar
usuário e senha para realizar as requisições. Após configurar corretamente, tudo
funcionou.

Com a ajuda da Laura, executamos os testes, identificamos um erro remanescente,
corrigimos e o \textbf{MR foi aprovado com sucesso}. Foi uma sprint marcante tanto
pelo contexto externo difícil quanto pelo aprendizado técnico: aprendi sobre
reutilização de entidades, autenticação em APIs e o processo completo de resolução
de problemas em ambiente Docker.
